\ifdefined\inMainDoc\else
\documentclass[12pt,a4paper]{article}
% ================================================================
%  PRÉAMBULE COMMUN - ANNALES CORRIGÉES
%  Université de Kara - FaST - Licence Fondamentale MPC
%  Design optimisé impression noir et blanc
% ================================================================

% -- ENCODAGE & LANGUE --
\usepackage[utf8]{inputenc}
\usepackage[T1]{fontenc}
\usepackage[french]{babel}
\usepackage{lmodern}
\usepackage{microtype}

% -- MISE EN PAGE --
\usepackage[a4paper, top=2.8cm, bottom=2.5cm, left=2.5cm, right=2.5cm, headheight=14pt]{geometry}
\usepackage{setspace}
\setstretch{1.2}
\usepackage{parskip}
\setlength{\parindent}{1.2em}
\setlength{\parskip}{4pt}

% -- MATHS --
\usepackage{amsmath, amssymb, mathtools, amsthm}
\usepackage{mathrsfs}
\usepackage{dsfont}
\usepackage{bm}
\usepackage{cancel}
\usepackage{textcomp}
% Définition de \oiint (intégrale de surface fermée) sans package supplémentaire
\newcommand{\oiint}{\mathop{\iint\!\!\!\!\!\!\!\bigcirc}\limits}

% -- PHYSIQUE & CHIMIE --
\usepackage{physics}
\usepackage{siunitx}
% Lorsque physics est chargé, siunitx omet \qty. On le restaure ici :
\AtBeginDocument{\RenewCommandCopy\qty\SI}
\sisetup{locale = FR, detect-all, output-decimal-marker={,}}
% Commande de substitution pour les unités posant problème avec pdflatex
% (ohm, degreeCelsius, micro, angstrom, percent utilisent des caractères non-ASCII)
\newcommand{\qtyohm}[1]{\ensuremath{\num{#1}\,\Omega}}
\newcommand{\qtydegc}[1]{\ensuremath{\num{#1}\,{}^\circ\text{C}}}
\newcommand{\qtyangstrom}[1]{\ensuremath{\num{#1}\,\text{\AA}}}
\newcommand{\qtypercent}[1]{\ensuremath{\num{#1}\,\%}}
\newcommand{\qtyum}[1]{\ensuremath{\num{#1}\,\mu\text{m}}}
\usepackage[version=4]{mhchem}

% -- GRAPHISMES --
\usepackage{graphicx}
\usepackage{float}
\usepackage{caption}
\usepackage{subcaption}
\usepackage{wrapfig}
\usepackage{tikz}
\usetikzlibrary{calc, arrows.meta, patterns, shapes.geometric}
\usepackage{pgfplots}
\pgfplotsset{compat=1.18}

% -- TABLEAUX --
\usepackage{booktabs}
\usepackage{array}
\usepackage{tabularx}
\usepackage{multirow}

% -- LISTES --
\usepackage{enumitem}
\setlist[enumerate]{topsep=3pt, itemsep=2pt, parsep=0pt}
\setlist[itemize]{topsep=3pt, itemsep=2pt, parsep=0pt, label={.}}

% -- BOÎTES (noir et blanc) --
\usepackage{mdframed}
\usepackage{tcolorbox}
\tcbuselibrary{skins, breakable, theorems}

% Boîte EXERCICE : encadré avec filet épais, fond blanc
\newmdenv[
  linewidth=1.2pt,
  linecolor=black,
  roundcorner=3pt,
  skipabove=10pt,
  skipbelow=6pt,
  innertopmargin=8pt,
  innerbottommargin=8pt,
  innerleftmargin=10pt,
  innerrightmargin=10pt,
  backgroundcolor=white,
  frametitle={\bfseries\small Exercice},
  frametitlealignment=\raggedright,
  frametitleaboveskip=4pt,
  frametitlebelowskip=4pt,
]{boiteexercice}

% Boîte CORRIGÉ : fond gris très clair + filet fin
\newmdenv[
  linewidth=0.6pt,
  linecolor=black,
  leftline=true,
  rightline=false,
  topline=false,
  bottomline=false,
  linecolor=black,
  innerleftmargin=12pt,
  innerrightmargin=8pt,
  innertopmargin=6pt,
  innerbottommargin=6pt,
  backgroundcolor=gray!8,
  skipabove=4pt,
  skipbelow=8pt,
]{boitecorrige}

% Boîte RÉSUMÉ DE COURS : double encadré
\newmdenv[
  linewidth=0.8pt,
  linecolor=black,
  roundcorner=2pt,
  backgroundcolor=gray!12,
  innertopmargin=10pt,
  innerbottommargin=10pt,
  innerleftmargin=12pt,
  innerrightmargin=12pt,
  skipabove=12pt,
  skipbelow=8pt,
]{boiteresume}

% Boîte ATTENTION/REMARQUE : barre gauche épaisse
\newmdenv[
  leftline=true,
  rightline=false,
  topline=false,
  bottomline=false,
  linewidth=3pt,
  linecolor=black,
  backgroundcolor=gray!5,
  innerleftmargin=12pt,
  innerrightmargin=8pt,
  innertopmargin=5pt,
  innerbottommargin=5pt,
  skipabove=6pt,
  skipbelow=6pt,
]{boiteremarque}

% -- DIFFICULTÉ DES EXERCICES --
\newcommand{\facile}{$\star$}
\newcommand{\moyen}{$\star\!\star$}
\newcommand{\difficile}{$\star\!\star\!\star$}
\newcommand{\niveau}[1]{\hfill\textbf{#1}\hspace{2pt}}

% -- EN-TÊTES ET PIEDS DE PAGE --
\usepackage{fancyhdr}
\pagestyle{fancy}
\fancyhf{}
\renewcommand{\headrulewidth}{0.5pt}
\renewcommand{\footrulewidth}{0.3pt}

% -- HYPERLIENS --
\usepackage[hidelinks]{hyperref}
\hypersetup{
    colorlinks=false,
    pdftitle={Annale Corrigée - Licence MPC - Université de Kara}
}

% -- TITRES --
\usepackage{titlesec}
\titleformat{\section}{\Large\bfseries}{\thesection.}{0.8em}{}[\vspace{2pt}\hrule\vspace{4pt}]
\titleformat{\subsection}{\large\bfseries}{\thesubsection.}{0.7em}{}
\titleformat{\subsubsection}{\normalsize\bfseries}{\thesubsubsection.}{0.6em}{}

% -- THÉORÈMES --
\theoremstyle{definition}
\newtheorem{defn}{Définition}[section]
\newtheorem{prop}{Propriété}[section]
\newtheorem{thm}{Théorème}[section]
\newtheorem{corol}{Corollaire}[section]
\newtheorem{exemple}{Exemple}[section]

% -- COMMANDES MATHÉMATIQUES UTILES --
\newcommand{\vect}[1]{\overrightarrow{#1}}
\newcommand{\R}{\mathbb{R}}
\newcommand{\N}{\mathbb{N}}
\newcommand{\Z}{\mathbb{Z}}
\newcommand{\C}{\mathbb{C}}
\renewcommand{\d}{\mathrm{d}}
\newcommand{\deriv}[2]{\dfrac{\d #1}{\d #2}}
\newcommand{\dpart}[2]{\dfrac{\partial #1}{\partial #2}}
\newcommand{\rot}{\overrightarrow{\mathrm{rot}}\,}
\renewcommand{\div}{\mathrm{div}\,}
\renewcommand{\grad}{\overrightarrow{\mathrm{grad}}\,}
\newcommand{\e}{\mathrm{e}}
\newcommand{\ii}{\mathrm{i}}
\renewcommand{\Tr}{\mathrm{Tr}}

% Numérotation des équations par section
\numberwithin{equation}{section}

% -- RÉSULTATS ENCADRÉS --
\newcommand{\res}[1]{\[\boxed{#1}\]}

% -- REMARQUE INLINE --
\newcommand{\rmq}[1]{\begin{boiteremarque}\textbf{Remarque.} #1\end{boiteremarque}}

% -- MISE EN VALEUR DU RÉSULTAT FINAL --
\newcommand{\reponse}[1]{%
  \begin{center}
    \fbox{\begin{minipage}{0.92\textwidth}\centering #1\end{minipage}}
  \end{center}%
}


\lhead{\small\textit{Calcul Différentiel - Licence 1}}
\rhead{\small\textit{FaST, Université de Kara}}
\cfoot{\thepage}

\begin{document}
\fi

\begin{center}
  {\LARGE\bfseries CALCUL DIFFÉRENTIEL DANS $\mathbb{R}$}\\[0.3cm]
  {\normalsize Licence 1 - Mathématiques, Physique et Chimie}\\[0.1cm]
  \rule{0.7\textwidth}{0.8pt}
\end{center}

\section*{Résumé du cours}

\begin{boiteresume}

\textbf{1. Limites.} $\lim_{x\to a}f(x)=L$ signifie que $f(x)$ se rapproche de $L$ quand $x$ s'approche de $a$. Règles : somme, produit, quotient (si le dénominateur ne s'annule pas). Formes indéterminées : $0/0$, $\infty/\infty$, $0\cdot\infty$, etc.

\textbf{2. Dérivation.} $f'(a) = \lim_{h\to0}\frac{f(a+h)-f(a)}{h}$ (si la limite existe). Règles de dérivation :
\begin{itemize}[nosep]
  \item $(u+v)' = u'+v'$, $(uv)' = u'v+uv'$, $(u/v)' = (u'v-uv')/v^2$
  \item $(f\circ g)'(x) = f'(g(x))\cdot g'(x)$ (dérivée de la composée)
  \item Dérivées usuelles : $(\e^x)'=\e^x$, $(\ln x)'=1/x$, $(\sin x)'=\cos x$, $(\cos x)'=-\sin x$, $(x^n)'=nx^{n-1}$
\end{itemize}

\textbf{3. Développements limités (DL).} En $a=0$ :
\[
\e^x = 1+x+\frac{x^2}{2!}+\cdots+\frac{x^n}{n!}+o(x^n)
\]
\[
\sin x = x - \frac{x^3}{6}+\cdots, \quad \cos x = 1-\frac{x^2}{2}+\cdots, \quad \ln(1+x) = x-\frac{x^2}{2}+\cdots
\]
\[
(1+x)^\alpha = 1+\alpha x+\frac{\alpha(\alpha-1)}{2}x^2+\cdots
\]

\textbf{4. Règle de L'Hôpital.} Si $\lim_{x\to a}f(x)=\lim_{x\to a}g(x)=0$ (ou $\pm\infty$) :
\[
\lim_{x\to a}\frac{f(x)}{g(x)} = \lim_{x\to a}\frac{f'(x)}{g'(x)}
\]

\textbf{5. Étude de fonctions.} Pour étudier $f$ : domaine de définition, parité, limites aux bornes, dérivée (signe et tableau de variations), extrema locaux et globaux, convexité ($f''$).

\textbf{6. Théorème de Rolle et des accroissements finis.} Si $f$ est continue sur $[a,b]$ et dérivable sur $]a,b[$, il existe $c\in]a,b[$ tel que $f(b)-f(a) = f'(c)(b-a)$.

\end{boiteresume}

\newpage
\section{Limites et continuité}

\subsection*{Exercice 1 - Calcul de limites \niveau{\facile}}

\begin{boiteexercice}
Calculer les limites suivantes :
\begin{enumerate}[label=(\alph*)]
  \item $\displaystyle\lim_{x\to 2}\frac{x^2-4}{x-2}$
  \item $\displaystyle\lim_{x\to 0}\frac{\sin(3x)}{x}$
  \item $\displaystyle\lim_{x\to+\infty}\frac{3x^2-2x+1}{x^2+5}$
  \item $\displaystyle\lim_{x\to 0}\frac{1-\cos x}{x^2}$
  \item $\displaystyle\lim_{x\to 0}x\ln x$
\end{enumerate}
\end{boiteexercice}

\begin{boitecorrige}
\textbf{(a)} Forme $0/0$. On factorise : $\frac{x^2-4}{x-2}=\frac{(x-2)(x+2)}{x-2}=x+2\to 4$.

\textbf{(b)} $\frac{\sin(3x)}{x}=3\cdot\frac{\sin(3x)}{3x}\to 3\times1 = 3$ (car $\lim_{u\to0}\frac{\sin u}{u}=1$).

\textbf{(c)} On divise par $x^2$ : $\frac{3-2/x+1/x^2}{1+5/x^2}\to\frac{3}{1}=3$.

\textbf{(d)} Par DL : $1-\cos x \approx \frac{x^2}{2}$ au voisinage de 0, donc $\frac{1-\cos x}{x^2}\to\frac{1}{2}$.

\textbf{(e)} Forme indéterminée $0\cdot(-\infty)$. On pose $t=1/x\to+\infty$ : $x\ln x = \frac{\ln(1/t)}{1/x}$... Plus simplement, avec L'Hôpital sur $\frac{\ln x}{1/x}$ :
$\lim_{x\to0^+}\frac{\ln x}{1/x}\stackrel{H}{=}\lim_{x\to0^+}\frac{1/x}{-1/x^2}=\lim_{x\to0^+}(-x)=0$.
Donc $\lim_{x\to0^+}x\ln x = 0$.
\end{boitecorrige}

\section{Dérivées et étude de fonctions}

\subsection*{Exercice 2 - Calcul de dérivées \niveau{\facile}}

\begin{boiteexercice}
Calculer la dérivée des fonctions suivantes :
\begin{enumerate}[label=(\alph*)]
  \item $f(x) = 3x^4 - 2x^2 + 5x - 1$
  \item $g(x) = \e^{x^2+1}$
  \item $h(x) = \sin(x)\ln(x)$
  \item $u(x) = \frac{x^2+1}{x-1}$
  \item $v(x) = \arctan(2x)$
  \item $w(x) = x^x$ (pour $x>0$)
\end{enumerate}
\end{boiteexercice}

\begin{boitecorrige}
\textbf{(a)} $f'(x) = 12x^3 - 4x + 5$.

\textbf{(b)} Composée : $g'(x) = 2x\,\e^{x^2+1}$.

\textbf{(c)} Produit : $h'(x) = \cos(x)\ln(x) + \sin(x)\cdot\frac{1}{x}$.

\textbf{(d)} Quotient : $u'(x) = \frac{2x(x-1)-(x^2+1)\cdot1}{(x-1)^2} = \frac{2x^2-2x-x^2-1}{(x-1)^2} = \frac{x^2-2x-1}{(x-1)^2}$.

\textbf{(e)} $(\arctan)' = \frac{1}{1+u^2}$, règle de composition : $v'(x) = \frac{2}{1+(2x)^2} = \frac{2}{1+4x^2}$.

\textbf{(f)} On écrit $x^x = \e^{x\ln x}$, puis :
$w'(x) = \e^{x\ln x}\cdot(x\ln x)' = x^x(\ln x + 1)$.
\end{boitecorrige}

\subsection*{Exercice 3 - Étude de fonction \niveau{\moyen}}

\begin{boiteexercice}
Étudier et tracer le graphe de $f(x) = x\,\e^{-x}$ :
\begin{enumerate}
  \item Domaine de définition et limites en $\pm\infty$.
  \item Calculer $f'(x)$ et dresser le tableau de variations.
  \item Déterminer les extrema et les inflexions.
  \item Tracer le graphe.
\end{enumerate}
\end{boiteexercice}

\begin{boitecorrige}
\textbf{1.} $\mathcal{D}_f = \mathbb{R}$. En $-\infty$ : $x\to-\infty$ et $\e^{-x}\to+\infty$, donc $f(x)\to-\infty$. En $+\infty$ : $f(x)=x/\e^x\to0^+$ (exponentielle l'emporte). Asymptote horizontale $y=0$ en $+\infty$.

\textbf{2.} $f'(x) = \e^{-x} + x(-\e^{-x}) = (1-x)\e^{-x}$.

$f'(x)>0$ si $x<1$, $f'(x)<0$ si $x>1$.

\begin{center}
\begin{tabular}{|c|cccccc|}
\hline
$x$ & $-\infty$ & & $1$ & & $+\infty$\\
\hline
$f'(x)$ & & $+$ & $0$ & $-$ &\\
\hline
$f(x)$ & $-\infty$ & $\nearrow$ & $\e^{-1}$ & $\searrow$ & $0$\\
\hline
\end{tabular}
\end{center}

\textbf{3.} Maximum en $x=1$ : $f(1)=\e^{-1}\approx0{,}368$.

$f''(x) = (-1)\e^{-x}+(1-x)(-\e^{-x}) = -(2-x)\e^{-x}=(x-2)\e^{-x}$.
$f''=0$ en $x=2$. Point d'inflexion en $(2,\;2\e^{-2})$.

\textbf{4.} La courbe part de $-\infty$ en croissant, atteint le maximum $\e^{-1}$ en $x=1$, puis décroît en tendant asymptotiquement vers $0$. Elle a un point d'inflexion en $(2, 2/\e^2)$.
\end{boitecorrige}

\section{Développements limités et applications}

\subsection*{Exercice 4 - Développements limités \niveau{\moyen}}

\begin{boiteexercice}
Écrire les développements limités suivants au voisinage de $x=0$, à l'ordre indiqué :
\begin{enumerate}[label=(\alph*)]
  \item $\frac{1}{1+x}$ à l'ordre 4
  \item $\sqrt{1+x}$ à l'ordre 3
  \item $\ln(1+\sin x)$ à l'ordre 3
  \item $\frac{\e^x-1}{\sin x}$ à l'ordre 2
\end{enumerate}
\end{boiteexercice}

\begin{boitecorrige}
\textbf{(a)} Série géométrique : $\frac{1}{1+x} = 1-x+x^2-x^3+x^4+o(x^4)$.

\textbf{(b)} Formule $(1+x)^\alpha$ avec $\alpha=1/2$ :
\begin{align*}
\sqrt{1+x} &= 1+\frac{1}{2}x+\frac{(1/2)(-1/2)}{2}x^2+\frac{(1/2)(-1/2)(-3/2)}{6}x^3+\cdots \\
&= 1+\frac{x}{2}-\frac{x^2}{8}+\frac{x^3}{16}+o(x^3)
\end{align*}

\textbf{(c)} On compose : $\sin x = x - x^3/6 + o(x^3)$, donc en posant $u = \sin x$ :
$\ln(1+\sin x) = \ln(1+u)$ avec $u=x-x^3/6+\cdots$
En utilisant $\ln(1+u)=u-u^2/2+u^3/3-\cdots$ avec $u=x-x^3/6+\cdots$ :
\[
= \left(x-\frac{x^3}{6}\right) - \frac{(x-x^3/6)^2}{2}+\frac{x^3}{3}+o(x^3)
= x - \frac{x^2}{2} - \frac{x^3}{6} + \frac{x^3}{3} + o(x^3)
= x - \frac{x^2}{2} + \frac{x^3}{6} + o(x^3)
\]

\textbf{(d)} $\e^x-1 = x+x^2/2+o(x^2)$, $\sin x = x+o(x^2)$ :
\begin{align*}
\frac{\e^x-1}{\sin x} &= \frac{x+x^2/2+o(x^2)}{x+o(x^2)} = \frac{1+x/2+o(x)}{1+o(x)} \\
&= \left(1+\frac{x}{2}+o(x)\right)(1+o(x)) = 1+\frac{x}{2}+o(x)
\end{align*}
\end{boitecorrige}

\subsection*{Exercice 5 - Règle de L'Hôpital \niveau{\moyen}}

\begin{boiteexercice}
Calculer les limites suivantes à l'aide de la règle de L'Hôpital :
\begin{enumerate}[label=(\alph*)]
  \item $\displaystyle\lim_{x\to0}\frac{\e^x-1-x}{x^2}$
  \item $\displaystyle\lim_{x\to0}\frac{\ln(1+x)-x+x^2/2}{x^3}$
  \item $\displaystyle\lim_{x\to+\infty}\frac{\ln x}{\sqrt{x}}$
  \item $\displaystyle\lim_{x\to1}\frac{x^3-3x+2}{x^3-x^2-x+1}$
\end{enumerate}
\end{boiteexercice}

\begin{boitecorrige}
\textbf{(a)} Forme $0/0$ :
$\frac{\e^x-1-x}{x^2}\stackrel{H}{=}\frac{\e^x-1}{2x}\stackrel{H}{=}\frac{\e^x}{2}\to\frac{1}{2}$.

\textbf{(b)} Forme $0/0$ :
$\frac{\ln(1+x)-x+x^2/2}{x^3}\stackrel{H}{=}\frac{\frac{1}{1+x}-1+x}{3x^2}=\frac{\frac{x^2}{1+x}}{3x^2}\cdot\frac{1}{...}$

Plus directement par DL : $\ln(1+x) = x-x^2/2+x^3/3+o(x^3)$, donc le numérateur vaut $x^3/3+o(x^3)$, et la limite est $\frac{1}{3}$.

\textbf{(c)} Forme $\infty/\infty$ :
$\frac{\ln x}{\sqrt{x}}\stackrel{H}{=}\frac{1/x}{1/(2\sqrt{x})}=\frac{2\sqrt{x}}{x}=\frac{2}{\sqrt{x}}\to0$.

\textbf{(d)} Forme $0/0$. Le numérateur : $x=1$ est racine de $x^3-3x+2$ (vérifié : $1-3+2=0$), donc $(x-1)$ est facteur ; $x^3-3x+2=(x-1)(x^2+x-2)=(x-1)^2(x+2)$.

Le dénominateur : $x=1$ est racine de $x^3-x^2-x+1$ (vérifié : $1-1-1+1=0$) ; $x^3-x^2-x+1=(x-1)(x^2-1)=(x-1)^2(x+1)$.
\[
\lim_{x\to1}\frac{(x-1)^2(x+2)}{(x-1)^2(x+1)} = \lim_{x\to1}\frac{x+2}{x+1} = \frac{3}{2}
\]
\end{boitecorrige}

\subsection*{Exercice 6 - Extrema et optimisation \niveau{\difficile}}

\begin{boiteexercice}
Un fermier dispose de \qty{120}{m} de clôture pour délimiter un enclos rectangulaire adossé à un mur (le mur constitue l'un des côtés, pas besoin de clôture de ce côté).
\begin{enumerate}
  \item Exprimer l'aire $A$ de l'enclos en fonction de la longueur $x$ du côté parallèle au mur.
  \item Déterminer $x$ qui maximise $A$.
  \item Quelle est l'aire maximale ?
\end{enumerate}
\end{boiteexercice}

\begin{boitecorrige}
\textbf{1.} Notons $x$ la longueur du côté parallèle au mur et $y$ la longueur des deux côtés perpendiculaires. La clôture utilise trois côtés : $x + 2y = 120$, donc $y = (120-x)/2 = 60-x/2$.

L'aire : $A(x) = x\cdot y = x(60-x/2) = 60x - x^2/2$, avec $0<x<120$.

\textbf{2.} $A'(x) = 60 - x = 0 \Rightarrow x = 60$.

$A''(x) = -1 < 0$ : c'est bien un maximum.

\textbf{3.} $A_{\max} = A(60) = 60\times60 - 60^2/2 = 3600 - 1800 = \qty{1800}{m^2}$.

Les dimensions optimales sont $x = \qty{60}{m}$ (côté parallèle au mur) et $y = \qty{30}{m}$ (côtés perpendiculaires).
\end{boitecorrige}

\section{Dérivées partielles et fonctions multivariables}

\subsection*{Exercice 7 - Dérivées partielles \niveau{\facile}}

\begin{boiteexercice}
Calculer toutes les dérivées partielles du premier et du second ordre des fonctions suivantes :
\begin{enumerate}[label=(\alph*)]
  \item $f(x,y) = x^2 y + xy^3 - 2x + 3y$
  \item $g(x,y) = \e^{x+2y}\sin(xy)$
  \item $h(x,y,z) = x^2+y^2+z^2$ (distance au carré à l'origine)
\end{enumerate}
\end{boiteexercice}

\begin{boitecorrige}
\textbf{(a)} $\frac{\partial f}{\partial x} = 2xy+y^3-2$, $\frac{\partial f}{\partial y} = x^2+3xy^2+3$, $\frac{\partial^2 f}{\partial x^2}=2y$, $\frac{\partial^2 f}{\partial y^2}=6xy$, $\frac{\partial^2 f}{\partial x\partial y}=2x+3y^2$ (symétrique de Schwarz).

\textbf{(b)} Par la règle du produit et de la composée :
\[
\frac{\partial g}{\partial x} = \e^{x+2y}\sin(xy) + \e^{x+2y}y\cos(xy) = \e^{x+2y}(\sin(xy)+y\cos(xy))
\]
\[
\frac{\partial g}{\partial y} = 2\e^{x+2y}\sin(xy) + \e^{x+2y}x\cos(xy) = \e^{x+2y}(2\sin(xy)+x\cos(xy))
\]

\textbf{(c)} $\frac{\partial h}{\partial x}=2x$, $\frac{\partial h}{\partial y}=2y$, $\frac{\partial h}{\partial z}=2z$.
$\frac{\partial^2 h}{\partial x^2} = 2$, toutes les dérivées croisées nulles. Le laplacien : $\Delta h = 2+2+2=6$.
\end{boitecorrige}

\subsection*{Exercice 8 - Différentielle totale et approximation \niveau{\moyen}}

\begin{boiteexercice}
\begin{enumerate}
  \item La pression $P$ d'un gaz parfait est donnée par $P = nRT/V$. Calculer la différentielle totale $\d P$.
  \item On mesure $T = \qty{300}{K}$ avec une erreur $\delta T = \qty{2}{K}$ et $V = \qty{2}{L}$ avec $\delta V = \qty{0.05}{L}$, pour $n = 1$ mol et $R = \qty{8.314}{J.mol^{-1}.K^{-1}}$. Estimer l'erreur relative sur $P$.
  \item Calculer la différentielle de $f(x,y) = \arctan(y/x)$ et en déduire une approximation de $f(1{,}01; 0{,}99)$ à partir de $f(1;1)$.
\end{enumerate}
\end{boiteexercice}

\begin{boitecorrige}
\textbf{1.} $P = nRT/V$, donc :
\[
\d P = \frac{\partial P}{\partial T}\d T + \frac{\partial P}{\partial V}\d V = \frac{nR}{V}\d T - \frac{nRT}{V^2}\d V
\]

\textbf{2.} Erreur relative : $\frac{|\delta P|}{P} \leq \frac{|\delta T|}{T} + \frac{|\delta V|}{V} = \frac{2}{300}+\frac{0{,}05}{2} = 0{,}00667+0{,}025 \approx 3{,}2\,\%$.

Valeur de $P$ : $P = 1\times8{,}314\times300/0{,}002 = \qty{1.247e6}{Pa}$. Erreur absolue $\delta P \approx 3{,}2\,\%\times P \approx \qty{4e4}{Pa}$.

\textbf{3.} La différentielle de $f = \arctan(y/x)$ s'écrit $\d f = \frac{\partial f}{\partial x}\d x + \frac{\partial f}{\partial y}\d y$. En appliquant $\frac{\d}{\d u}\arctan u = \frac{1}{1+u^2}$ et la règle de la chaîne :
\[
\frac{\partial f}{\partial x} = \frac{1}{1+(y/x)^2}\cdot\left(-\frac{y}{x^2}\right) = \frac{-y}{x^2+y^2}
\]
\[
\frac{\partial f}{\partial y} = \frac{1}{1+(y/x)^2}\cdot\frac{1}{x} = \frac{x}{x^2+y^2}
\]
En $(1,1)$ : $f(1,1)=\arctan(1)=\pi/4$. $\delta x = 0{,}01$, $\delta y = -0{,}01$.
\[
f(1{,}01;0{,}99) \approx f(1,1) + \frac{-1}{2}\times0{,}01+\frac{1}{2}\times(-0{,}01) = \frac{\pi}{4}-\frac{0{,}01}{2}-\frac{0{,}01}{2} = \frac{\pi}{4}-0{,}01 \approx 0{,}775
\]
\end{boitecorrige}

\subsection*{Exercice 9 - Théorème des accroissements finis \niveau{\moyen}}

\begin{boiteexercice}
\begin{enumerate}
  \item Appliquer le théorème de Rolle à $f(x) = x^3 - 3x$ sur $[-\sqrt{3}, \sqrt{3}]$.
  \item Appliquer le TAF à $g(x) = x^3$ sur $[0, 2]$ et trouver $c$.
  \item Montrer que $|\sin x - \sin y| \leq |x-y|$ pour tout $x, y \in \mathbb{R}$.
  \item Montrer que $\sqrt{1+x} < 1+x/2$ pour $x > 0$.
\end{enumerate}
\end{boiteexercice}

\begin{boitecorrige}
\textbf{1.} $f(-\sqrt{3}) = -3\sqrt{3}+3\sqrt{3} = 0 = f(\sqrt{3})$. Les hypothèses du théorème de Rolle sont vérifiées ($f$ continue et dérivable). $f'(x)=3x^2-3=0$ donne $x=\pm1\in]-\sqrt{3},\sqrt{3}[$. Il existe bien un point où $f'=0$.

\textbf{2.} $g(2)-g(0) = 8$. TAF : $8 = g'(c)\times(2-0) = 3c^2\times2$. Donc $c^2 = 4/3$, $c = 2/\sqrt{3} = 2\sqrt{3}/3 \approx 1{,}155 \in ]0,2[$.

\textbf{3.} Appliquons le TAF à $\sin$ sur $[y, x]$ (si $x>y$) : il existe $c\in]y,x[$ tel que
$|\sin x-\sin y| = |\cos c|\cdot|x-y| \leq |x-y|$, car $|\cos c|\leq 1$.

\textbf{4.} Posons $h(x) = 1+x/2 - \sqrt{1+x}$ pour $x>0$. $h(0)=0$.
$h'(x) = 1/2 - \frac{1}{2\sqrt{1+x}} = \frac{1}{2}\left(1-\frac{1}{\sqrt{1+x}}\right)$.
Pour $x>0$, $\sqrt{1+x}>1$, donc $1/\sqrt{1+x}<1$, et $h'(x)>0$. Ainsi $h$ est croissante et $h(x)>h(0)=0$ pour $x>0$.
\end{boitecorrige}

\subsection*{Exercice 10 - Extrema de fonctions à deux variables \niveau{\difficile}}

\begin{boiteexercice}
Trouver et classifier les points critiques de :
\begin{enumerate}[label=(\alph*)]
  \item $f(x,y) = x^2 + y^2 - 2x - 4y + 1$
  \item $g(x,y) = x^3 + y^3 - 3xy$
  \item $h(x,y) = x^2 - y^2$
\end{enumerate}
\end{boiteexercice}

\begin{boitecorrige}
On utilise le critère de la matrice hessienne $H = \begin{pmatrix}f_{xx}&f_{xy}\\f_{xy}&f_{yy}\end{pmatrix}$.

\textbf{(a)} $f_x = 2x-2=0 \Rightarrow x=1$; $f_y = 2y-4=0 \Rightarrow y=2$. Point critique $(1,2)$.
$H = \begin{pmatrix}2&0\\0&2\end{pmatrix}$, $\det H = 4>0$, $f_{xx}=2>0$. \textbf{Minimum} en $(1,2)$, $f(1,2) = 1+4-2-8+1=-4$.

\textbf{(b)} $g_x = 3x^2-3y=0$, $g_y = 3y^2-3x=0$. De la 1re : $y=x^2$. Substitution : $3x^4-3x=0$, soit $3x(x^3-1)=0$. Solutions : $(0,0)$ et $(1,1)$.

En $(0,0)$ : $H = \begin{pmatrix}0&-3\\-3&0\end{pmatrix}$, $\det H = -9<0$ : \textbf{point selle}.

En $(1,1)$ : $H = \begin{pmatrix}6&-3\\-3&6\end{pmatrix}$, $\det H = 36-9=27>0$, $f_{xx}=6>0$ : \textbf{minimum} en $(1,1)$, $g(1,1)=1+1-3=-1$.

\textbf{(c)} $h_x = 2x=0$, $h_y = -2y=0$. Point critique unique $(0,0)$.
$H = \begin{pmatrix}2&0\\0&-2\end{pmatrix}$, $\det H = -4<0$ : \textbf{point selle} (la surface est un paraboloïde hyperbolique).
\end{boitecorrige}

\subsection*{Exercice 11 - Optimisation sous contrainte (Lagrange) \niveau{\difficile}}

\begin{boiteexercice}
\begin{enumerate}
  \item Trouver le rectangle d'aire maximale inscrit dans une ellipse $\frac{x^2}{a^2}+\frac{y^2}{b^2}=1$, en utilisant la méthode des multiplicateurs de Lagrange.
  \item Trouver la distance minimale de l'origine au plan $2x+y-3z=6$.
\end{enumerate}
\end{boiteexercice}

\begin{boitecorrige}
\textbf{1.} On maximise l'aire $A = 4xy$ (un quart du rectangle, dans le premier quadrant) sous la contrainte $g(x,y) = x^2/a^2+y^2/b^2-1=0$.

Conditions de Lagrange : $\overrightarrow{\text{grad}}\,A = \lambda\,\overrightarrow{\text{grad}}\,g$ :
\[
4y = \lambda\cdot\frac{2x}{a^2} \qquad 4x = \lambda\cdot\frac{2y}{b^2}
\]
Division membre à membre : $y/x = x b^2/(y a^2)$, soit $y^2 a^2 = x^2 b^2$, donc $y = bx/a$.

Substitution dans la contrainte : $x^2/a^2 + b^2x^2/(a^2b^2) = 2x^2/a^2 = 1$, $x = a/\sqrt{2}$, $y = b/\sqrt{2}$.

Aire maximale : $A = 4\times\frac{a}{\sqrt{2}}\times\frac{b}{\sqrt{2}} = 2ab$.

\textbf{2.} On minimise $d^2 = x^2+y^2+z^2$ sous $2x+y-3z=6$. Conditions de Lagrange :
\[
2x = 2\lambda,\quad 2y=\lambda,\quad 2z=-3\lambda
\]
Soit $x=\lambda$, $y=\lambda/2$, $z=-3\lambda/2$. Substitution :
$2\lambda+\lambda/2-3\times(-3\lambda/2) = 2\lambda+\lambda/2+9\lambda/2 = 2\lambda+5\lambda = 7\lambda = 6$. Donc $\lambda = 6/7$.
\[
(x,y,z) = \left(\frac{6}{7},\frac{3}{7},-\frac{9}{7}\right), \quad d = \sqrt{\frac{36+9+81}{49}} = \frac{\sqrt{126}}{7} = \frac{3\sqrt{14}}{7} \approx 1{,}604
\]
\end{boitecorrige}

\subsection*{Exercice 12 - Série de Taylor et applications \niveau{\moyen}}

\begin{boiteexercice}
\begin{enumerate}
  \item Écrire le développement de Taylor de $f(x)=\cos x$ en $x=0$ jusqu'à l'ordre 6.
  \item Utiliser ce DL pour calculer une valeur approchée de $\cos(0{,}1)$ à $10^{-5}$ près.
  \item Montrer que pour $x$ proche de $\pi/2$, $\cos x \approx -(x-\pi/2)$ (DL d'ordre 1 en $a=\pi/2$).
  \item Calculer $\lim_{x\to0}\frac{1-\cos x}{x\sin x}$ en utilisant les DL.
\end{enumerate}
\end{boiteexercice}

\begin{boitecorrige}
\textbf{1.} $\cos x = 1-\frac{x^2}{2!}+\frac{x^4}{4!}-\frac{x^6}{6!}+o(x^6) = 1-\frac{x^2}{2}+\frac{x^4}{24}-\frac{x^6}{720}+o(x^6)$.

\textbf{2.} Pour $x=0{,}1$ :
\[
\cos(0{,}1) \approx 1-\frac{0{,}01}{2}+\frac{0{,}0001}{24} = 1-0{,}005+0{,}0000041\overline{6} \approx 0{,}995004
\]
Le terme d'ordre 6 : $\frac{(0{,}1)^6}{720} \approx \frac{10^{-6}}{720} < 10^{-5}$. L'approximation est correcte à $10^{-5}$ près.

\textbf{3.} En $a=\pi/2$, posons $u=x-\pi/2$. $\cos(\pi/2+u) = -\sin u$. Pour $u$ petit : $-\sin u\approx -u = -(x-\pi/2)$.

\textbf{4.} Par DL au voisinage de 0 :
\[
1-\cos x \approx \frac{x^2}{2}, \quad x\sin x \approx x\cdot x = x^2
\]
\[
\lim_{x\to0}\frac{1-\cos x}{x\sin x} = \lim_{x\to0}\frac{x^2/2}{x^2} = \frac{1}{2}
\]
\end{boitecorrige}


% ================================================================
\subsection*{Exercice 13 - Développement limité de $\sqrt{1+x}$ \niveau{\facile}}
% ================================================================
\begin{boiteexercice}
\begin{enumerate}
  \item En utilisant la formule du développement limité de $(1+u)^\alpha$ à l'ordre 3 en $u=0$ :
  \[
  (1+u)^\alpha = 1 + \alpha u + \frac{\alpha(\alpha-1)}{2}u^2 + \frac{\alpha(\alpha-1)(\alpha-2)}{6}u^3 + o(u^3)
  \]
  Écrire le développement limité de $f(x) = \sqrt{1+x}$ à l'ordre 3 en $x=0$.
  \item En déduire une valeur approchée de $\sqrt{1{,}04}$ avec une erreur estimée.
  \item Calculer de même le DL à l'ordre 2 de $\dfrac{1}{\sqrt{1+x}}$.
\end{enumerate}
\end{boiteexercice}

\begin{boitecorrige}
\textbf{1.} On applique avec $\alpha = 1/2$ et $u = x$ :
\begin{align*}
\alpha &= \tfrac{1}{2} \\
\frac{\alpha(\alpha-1)}{2} &= \frac{\tfrac{1}{2}\cdot(-\tfrac{1}{2})}{2} = -\frac{1}{8} \\
\frac{\alpha(\alpha-1)(\alpha-2)}{6} &= \frac{\tfrac{1}{2}\cdot(-\tfrac{1}{2})\cdot(-\tfrac{3}{2})}{6} = \frac{3/8}{6} = \frac{1}{16}
\end{align*}
\[
\sqrt{1+x} = 1 + \frac{x}{2} - \frac{x^2}{8} + \frac{x^3}{16} + o(x^3)
\]

\textbf{2.} Pour $x = 0{,}04$ :
\[
\sqrt{1{,}04} \approx 1 + \frac{0{,}04}{2} - \frac{(0{,}04)^2}{8} + \frac{(0{,}04)^3}{16}
= 1 + 0{,}02 - 0{,}0002 + 0{,}000004 \approx 1{,}0198
\]
Le terme d'ordre suivant est de l'ordre de $(0{,}04)^4/128 < 10^{-7}$. L'approximation est très précise (valeur réelle : $1{,}019804\ldots$).

\textbf{3.} Pour $\alpha = -1/2$ :
\[
\frac{1}{\sqrt{1+x}} = (1+x)^{-1/2} = 1 - \frac{x}{2} + \frac{3x^2}{8} + o(x^2)
\]
\end{boitecorrige}

% ================================================================
\subsection*{Exercice 14 - Étude complète de $f(x)=\dfrac{x^2-1}{x^2+1}$ \niveau{\moyen}}
% ================================================================
\begin{boiteexercice}
Soit $f(x) = \dfrac{x^2-1}{x^2+1}$.

\begin{enumerate}
  \item Déterminer le domaine de définition $\mathcal{D}_f$.
  \item Étudier la parité de $f$.
  \item Calculer les limites aux bornes du domaine.
  \item Calculer $f'(x)$ et en déduire le tableau de variation.
  \item Déterminer les asymptotes (horizontale, oblique, verticale).
  \item Tracer l'allure de la courbe.
\end{enumerate}
\end{boiteexercice}

\begin{boitecorrige}
\textbf{1. Domaine :} $x^2+1 > 0$ pour tout $x\in\mathbb{R}$, donc $\mathcal{D}_f = \mathbb{R}$.

\textbf{2. Parité :} $f(-x) = \dfrac{(-x)^2-1}{(-x)^2+1} = \dfrac{x^2-1}{x^2+1} = f(x)$ : $f$ est \textbf{paire}. La courbe est symétrique par rapport à l'axe $Oy$.

\textbf{3. Limites :}
\[
\lim_{x\to\pm\infty} f(x) = \lim_{x\to\pm\infty}\frac{x^2(1-1/x^2)}{x^2(1+1/x^2)} = 1
\]
$f(0) = \dfrac{-1}{1} = -1$.

\textbf{4. Dérivée :}
\[
f'(x) = \frac{2x(x^2+1) - (x^2-1)\cdot 2x}{(x^2+1)^2} = \frac{2x\bigl[(x^2+1)-(x^2-1)\bigr]}{(x^2+1)^2} = \frac{4x}{(x^2+1)^2}
\]
$f'(x) = 0 \Leftrightarrow x=0$ ; $f'(x)<0$ pour $x<0$ ; $f'(x)>0$ pour $x>0$.

\begin{center}
\begin{tabular}{c|ccccc}
$x$ & $-\infty$ & & $0$ & & $+\infty$ \\
\hline
$f'(x)$ & & $-$ & $0$ & $+$ & \\
\hline
$f(x)$ & $1$ & $\searrow$ & $-1$ & $\nearrow$ & $1$ \\
\end{tabular}
\end{center}

$f(0) = -1$ est un \textbf{minimum global}.

\textbf{5. Asymptotes :} $y=1$ est une \textbf{asymptote horizontale} en $\pm\infty$ (la courbe ne coupe jamais $y=1$ car $x^2-1=x^2+1$ est impossible). Pas d'asymptote verticale.

\textbf{6. Points remarquables :} $f(x)=0 \Leftrightarrow x^2=1 \Leftrightarrow x=\pm1$. Ordonnée à l'origine : $(0,-1)$.
\end{boitecorrige}

% ================================================================
\subsection*{Exercice 15 - Règle de l'Hôpital \niveau{\moyen}}
% ================================================================
\begin{boiteexercice}
Calculer les limites suivantes en utilisant la règle de l'Hôpital :

\begin{enumerate}
  \item $\displaystyle\lim_{x\to0}\frac{\sin x - x}{x^3}$
  \item $\displaystyle\lim_{x\to+\infty} x\ln\!\left(1+\frac{1}{x}\right)$
  \item $\displaystyle\lim_{x\to0}\frac{e^x - 1 - x}{x^2}$
\end{enumerate}
\end{boiteexercice}

\begin{boitecorrige}
\textbf{1.} Forme $0/0$. Application de l'Hôpital (3 fois) :
\[
\lim_{x\to0}\frac{\sin x - x}{x^3} \overset{H}{=} \lim_{x\to0}\frac{\cos x - 1}{3x^2} \overset{H}{=} \lim_{x\to0}\frac{-\sin x}{6x} \overset{H}{=} \lim_{x\to0}\frac{-\cos x}{6} = \boxed{-\frac{1}{6}}
\]

\textbf{Vérification par DL :} $\sin x = x - x^3/6 + o(x^3)$, donc $(\sin x - x)/x^3 \to -1/6$.

\textbf{2.} Forme $+\infty \cdot 0$. On pose $u = 1/x \to 0^+$ :
\[
x\ln\!\left(1+\frac{1}{x}\right) = \frac{\ln(1+u)}{u} \overset{H}{=} \frac{1/(1+u)}{1}\xrightarrow{u\to0} \boxed{1}
\]
Interprétation : $\ln(1+u) \sim u$ en $0$, donc $\ln(1+1/x) \sim 1/x$ quand $x\to+\infty$.

\textbf{3.} Forme $0/0$. Application de l'Hôpital (2 fois) :
\[
\lim_{x\to0}\frac{e^x-1-x}{x^2} \overset{H}{=} \lim_{x\to0}\frac{e^x-1}{2x} \overset{H}{=} \lim_{x\to0}\frac{e^x}{2} = \boxed{\frac{1}{2}}
\]
\end{boitecorrige}

% ================================================================
\subsection*{Exercice 16 - Dérivées partielles et théorème de Schwarz \niveau{\moyen}}
% ================================================================
\begin{boiteexercice}
Soit $f(x,y) = x^2y + \sin(xy)$.

\begin{enumerate}
  \item Calculer les dérivées partielles du premier ordre $\dfrac{\partial f}{\partial x}$ et $\dfrac{\partial f}{\partial y}$.
  \item Calculer les dérivées partielles du second ordre $\dfrac{\partial^2 f}{\partial x^2}$, $\dfrac{\partial^2 f}{\partial y^2}$, $\dfrac{\partial^2 f}{\partial x\partial y}$ et $\dfrac{\partial^2 f}{\partial y\partial x}$.
  \item Vérifier le théorème de Schwarz (symétrie des dérivées croisées).
  \item Calculer $\dfrac{\partial f}{\partial x}$ et $\dfrac{\partial f}{\partial y}$ au point $(0, \pi)$.
\end{enumerate}
\end{boiteexercice}

\begin{boitecorrige}
\textbf{1. Dérivées partielles du premier ordre :}
\[
\frac{\partial f}{\partial x} = 2xy + y\cos(xy)
\]
\[
\frac{\partial f}{\partial y} = x^2 + x\cos(xy)
\]

\textbf{2. Dérivées partielles du second ordre :}
\[
\frac{\partial^2 f}{\partial x^2} = \frac{\partial}{\partial x}(2xy + y\cos(xy)) = 2y + y\cdot(-y\sin(xy)) = 2y - y^2\sin(xy)
\]
\[
\frac{\partial^2 f}{\partial y^2} = \frac{\partial}{\partial y}(x^2 + x\cos(xy)) = x\cdot(-x\sin(xy)) = -x^2\sin(xy)
\]
\[
\frac{\partial^2 f}{\partial x\partial y} = \frac{\partial}{\partial x}(x^2 + x\cos(xy)) = 2x + \cos(xy) + x\cdot(-y\sin(xy)) = 2x + \cos(xy) - xy\sin(xy)
\]
\[
\frac{\partial^2 f}{\partial y\partial x} = \frac{\partial}{\partial y}(2xy + y\cos(xy)) = 2x + \cos(xy) + y\cdot(-x\sin(xy)) = 2x + \cos(xy) - xy\sin(xy)
\]

\textbf{3. Théorème de Schwarz :}
\[
\frac{\partial^2 f}{\partial x\partial y} = \frac{\partial^2 f}{\partial y\partial x} = 2x + \cos(xy) - xy\sin(xy) \;
\]
Les dérivées croisées sont égales car $f$ est de classe $\mathcal{C}^2$ sur $\mathbb{R}^2$.

\textbf{4. Au point $(0, \pi)$ :}
\[
\frac{\partial f}{\partial x}(0,\pi) = 2\cdot0\cdot\pi + \pi\cos(0) = \pi
\]
\[
\frac{\partial f}{\partial y}(0,\pi) = 0^2 + 0\cdot\cos(0) = 0
\]
\end{boitecorrige}

\ifdefined\inMainDoc\else\end{document}\fi
